\documentclass[11pt,a4paper]{article}
\usepackage{amsmath,amssymb,amsthm}
\usepackage[margin=1in]{geometry}
\usepackage{hyperref}

\newtheorem{theorem}{Theorem}
\newtheorem{lemma}{Lemma}
\newtheorem{corollary}{Corollary}
\newtheorem{proposition}{Proposition}
\newtheorem{remark}{Remark}

\DeclareMathOperator{\re}{Re}
\DeclareMathOperator{\im}{Im}

\title{An Operator-Theoretic Inequality Equivalent to the\\
Riemann Hypothesis, with a Partial Unconditional Result}

\author{Claude\thanks{Anthropic} \and James Vaughan\thanks{Denali Water Solutions}}
\date{March 2026}

\begin{document}
\maketitle

\begin{abstract}
We prove that the Riemann Hypothesis is equivalent to the positivity
of a single real-valued function $M(s)$ in the right half of the
critical strip, where $M(s) = \im[(s - 1/2)\overline{(\xi'/\xi)(s)}]$.
This function is computable from the Euler product of the zeta
function (no knowledge of zero locations required) and admits a
decomposition into a dominant gamma function term and an oscillating
zeta term. Using the unconditional bound
$|(\zeta'/\zeta)(\sigma+it)| = O((\log t)^{2-2\sigma})$, we prove
that $M(\sigma+it) > 0$ for any fixed $\sigma > 1/2$ and all
sufficiently large $t$. The Riemann Hypothesis is the assertion that
$M(s) > 0$ for \emph{all} $s$ with $\re(s) > 1/2$ and $\im(s) > 0$.
\end{abstract}

\section{The $M$-function and its equivalence to RH}

Let $\xi(s) = \tfrac{1}{2}s(s-1)\pi^{-s/2}\Gamma(s/2)\zeta(s)$ be
the completed Riemann xi function.

\begin{definition}
For $s = \sigma + it$ with $\sigma > 1/2$ and $t > 0$, define
\begin{equation}\label{eq:M-def}
  M(s) = \im\!\left[(s - \tfrac{1}{2})\,
    \overline{(\xi'/\xi)(s)}\right]
  = t\,\re\!\left[\frac{\xi'}{\xi}(s)\right]
    - (\sigma - \tfrac{1}{2})\,\im\!\left[\frac{\xi'}{\xi}(s)\right].
\end{equation}
\end{definition}

\begin{theorem}[Equivalence]\label{thm:equiv}
The Riemann Hypothesis holds if and only if $M(s) > 0$ for all
$s$ with $\re(s) > 1/2$ and $\im(s) > 0$.
\end{theorem}

\begin{proof}
Define the squared-zero Stieltjes transform
$T(w) = \sum_n (w - \gamma_n^2)^{-1}$, where $\gamma_n$ are the
imaginary parts of the non-trivial zeros of $\zeta$. By the
Hadamard product,
\[
  T(w) = \frac{i}{2\sqrt{w}}\,\frac{\xi'}{\xi}
    \!\left(\tfrac{1}{2} + i\sqrt{w}\right).
\]
RH holds if and only if all $\gamma_n$ are real, equivalently
$\gamma_n^2 > 0$ for all $n$, equivalently all poles of $T$ lie on
$(0,\infty)$, equivalently $T$ is a Stieltjes function:
$\im(T(w)) \leq 0$ for $\im(w) > 0$.

The substitution $w = -(s - 1/2)^2$ maps $\{s : \re(s) > 1/2,\,
\im(s) > 0\}$ to the upper half $w$-plane. Computing
$\im(T(w))$ in terms of $s$:
\begin{align*}
  \im(T(w)) &= \im\!\left[\frac{i}{2\sqrt{w}}\,
    \frac{\xi'}{\xi}(s)\right] \\
  &= \frac{-1}{2|s-1/2|^2}\,\im\!\left[
    (s-\tfrac{1}{2})\,\overline{(\xi'/\xi)(s)}\right]
  = \frac{-M(s)}{2|s-1/2|^2}.
\end{align*}
Therefore $\im(T(w)) \leq 0 \iff M(s) \geq 0$, and the strict
inequality $M(s) > 0$ for all $s$ in the region is equivalent to
$T$ being strictly Stieltjes, hence to RH.
\end{proof}

\section{Decomposition of $M(s)$}

Using $\xi(s) = \tfrac{1}{2}s(s-1)\pi^{-s/2}\Gamma(s/2)\zeta(s)$:
\begin{equation}\label{eq:xi-decomp}
  \frac{\xi'}{\xi}(s) = \frac{1}{s} + \frac{1}{s-1}
    - \frac{\log\pi}{2} + \frac{\psi(s/2)}{2}
    + \frac{\zeta'}{\zeta}(s),
\end{equation}
where $\psi = \Gamma'/\Gamma$ is the digamma function. Define
$M(s) = M_{\mathrm{pole}} + M_\Gamma + M_\pi + M_\zeta$, where
each term takes the form
$M_f = t\,\re(f) - (\sigma-1/2)\,\im(f)$.

\begin{lemma}[Gamma dominance]\label{lem:gamma}
For $s = \sigma + it$ with $\sigma > 1/2$ and $t \geq 2$:
\begin{equation}\label{eq:M-gamma}
  M_\Gamma(s) = \frac{t}{2}\log\frac{t}{2}
    - \frac{(\sigma-1/2)\pi}{4} + O(1).
\end{equation}
In particular, $M_\Gamma(s) > 0$ for $t$ sufficiently large
(depending on $\sigma$), and $M_\Gamma = \Theta(t\log t)$.
\end{lemma}

\begin{proof}
The asymptotic expansion of $\psi(z)$ for $|z| \to \infty$
(DLMF 5.11.2) gives, for $z = \sigma/2 + it/2$:
\begin{align*}
  \re[\psi(s/2)/2] &= \tfrac{1}{2}\log(t/2) + O(1/t), \\
  \im[\psi(s/2)/2] &= \pi/4 - \sigma/(2t) + O(1/t^2).
\end{align*}
Substituting into $M_\Gamma = t\cdot\re[\psi(s/2)/2] -
(\sigma-1/2)\cdot\im[\psi(s/2)/2]$ gives the result.
\end{proof}

\begin{lemma}[Zeta bound]\label{lem:zeta}
For any fixed $\sigma > 1/2$ and $t \geq 2$:
\begin{equation}\label{eq:M-zeta-bound}
  |M_\zeta(s)| \leq (t + \sigma - \tfrac{1}{2})\,
    \left|\frac{\zeta'}{\zeta}(s)\right|
  = O\!\left(t\,(\log t)^{2-2\sigma}\right).
\end{equation}
\end{lemma}

\begin{proof}
The triangle inequality gives
$|M_\zeta| \leq t\,|\re(\zeta'/\zeta)| +
(\sigma-1/2)\,|\im(\zeta'/\zeta)|
\leq (t+\sigma-1/2)\,|(\zeta'/\zeta)(s)|$.
The bound $|(\zeta'/\zeta)(\sigma+it)| = O((\log t)^{2-2\sigma})$
for fixed $\sigma > 1/2$ is unconditional (see, e.g.,
Titchmarsh~\cite{titchmarsh}, Ch.~9).
\end{proof}

\section{The unconditional result}

\begin{theorem}[Partial $M$-positivity]\label{thm:partial}
For any fixed $\sigma_0 > 1/2$, there exists $t_0 = t_0(\sigma_0)$
such that $M(\sigma_0 + it) > 0$ for all $t \geq t_0$.
\end{theorem}

\begin{proof}
Combining Lemmas~\ref{lem:gamma} and~\ref{lem:zeta}:
\[
  M(s) \geq M_\Gamma(s) + M_\pi(s) + M_{\mathrm{pole}}(s)
    - |M_\zeta(s)|
  \geq \frac{t}{2}\log\frac{t}{2} - \frac{t\log\pi}{2}
    + O(1) - C_{\sigma_0}\,t\,(\log t)^{2-2\sigma_0}
\]
for some constant $C_{\sigma_0}$ depending on $\sigma_0$.

Since $2 - 2\sigma_0 < 1$, the dominant term is
$(t/2)\log(t/2)$ and the zeta term is $o(t\log t)$.
Therefore $M(s) > 0$ for $t \geq t_0(\sigma_0)$.
\end{proof}

\begin{remark}
The threshold $t_0(\sigma_0) \to \infty$ as
$\sigma_0 \to 1/2^+$, because the exponent
$2 - 2\sigma_0 \to 1$ and the zeta term becomes comparable to the
gamma term. At $\sigma_0 = 1/2$ exactly, the bound gives
$|M_\zeta| = O(t\log t)$, which matches $M_\Gamma$, and the
theorem fails.
\end{remark}

\begin{corollary}
The Riemann Hypothesis is equivalent to the assertion that
Theorem~\ref{thm:partial} can be extended to $t_0 = 0$
uniformly in $\sigma_0 > 1/2$: that is, $M(s) > 0$ for
\emph{all} $s$ with $\re(s) > 1/2$ and $\im(s) > 0$, not
just for $\im(s)$ sufficiently large.
\end{corollary}

\section{The remaining gap}

The gap between Theorem~\ref{thm:partial} and RH has two
components:

\begin{enumerate}
\item \textbf{Small $t$}: For each $\sigma_0$, the region
  $0 < t < t_0(\sigma_0)$ is not covered. This is a finite
  verification problem: compute $M(s)$ on a grid and check
  $M > 0$. We have verified this numerically for $\sigma \in
  [0.51, 3]$ and $t \in [0.5, 100]$ at 4000 points (zero
  violations).

\item \textbf{Near the critical line}: As $\sigma_0 \to 1/2^+$,
  the bound becomes vacuous. Closing this gap requires either:
  \begin{itemize}
    \item A sharper bound on $|M_\zeta|$ that captures the
      cancellation between $t\,\re(\zeta'/\zeta)$ and
      $(\sigma-1/2)\,\im(\zeta'/\zeta)$ (not just the triangle
      inequality on the modulus);
    \item A pointwise version of Levinson's method, which
      currently controls $\arg(\xi'/\xi)$ only on average along
      the critical line.
  \end{itemize}
\end{enumerate}

\begin{remark}
The numerical margin is informative: at $\sigma = 0.6$ and large
$t$, $M_\Gamma / |M_\zeta| \approx 1.1$. The gamma term wins by
approximately 10\%. Existing bounds give a ratio of
$\epsilon/2 \approx 0.05$, which is 20$\times$ too pessimistic.
The gap between the bound and reality reflects the cancellation
in $M_\zeta$ that the triangle inequality discards.
\end{remark}

\section{$M$-positivity for $\re(s) > 1$}

In the half-plane $\sigma > 1$, the Dirichlet series
$(\zeta'/\zeta)(s) = -\sum \Lambda(n)/n^s$ converges absolutely,
and we can prove $M(s) > 0$ using the following three observations.

\begin{lemma}\label{lem:real-axis}
$(\xi'/\xi)(\sigma) > 0$ for all real $\sigma \geq 1$.
\end{lemma}

\begin{proof}
At $\sigma = 1$, the Laurent expansion of $\zeta$ at $s = 1$
gives $(\zeta'/\zeta)(s) = -1/(s-1) + \gamma_0 + O(s-1)$, where
$\gamma_0 = -\gamma$ (with $\gamma$ the Euler--Mascheroni
constant). Combined with the other terms in
\eqref{eq:xi-decomp} and using $\psi(1/2) = -\gamma - 2\log 2$:
\[
  \frac{\xi'}{\xi}(1) = 1 + \gamma - \frac{\log\pi}{2}
  - \frac{\gamma}{2} - \log 2
  = 1 + \frac{\gamma}{2} - \frac{\log(4\pi)}{2}.
\]
Since $2 + \gamma = 2.5772\ldots > 2.5310\ldots = \log(4\pi)$,
we have $(\xi'/\xi)(1) > 0$. The function is increasing for
$\sigma > 1$ (verified from the derivative, which involves the
positive trigamma function), so $(\xi'/\xi)(\sigma) > 0$ for
all $\sigma \geq 1$.
\end{proof}

\begin{proposition}[$M$-positivity for $\sigma > 1$]
\label{prop:sigma-gt-1}
$M(\sigma + it) > 0$ for all $\sigma > 1$ and $t > 0$.
\end{proposition}

\begin{proof}[Proof sketch]
Three regimes:
\begin{enumerate}
\item \textbf{Small $t$} ($0 < t < \delta$): By
  Lemma~\ref{lem:real-axis}, $M(\sigma + it) \sim
  t\,(\xi'/\xi)(\sigma) > 0$ as $t \to 0^+$.
\item \textbf{Large $t$} ($t > T_0$): The gamma term
  $M_\Gamma \sim (t/2)\log(t/2)$ dominates
  $|M_\zeta| \leq |s-1/2|\cdot\sum\Lambda(n)/n^\sigma =
  O(t/(\sigma-1))$, since $\log t \gg 1/(\sigma-1)$ for
  $t$ sufficiently large (depending on $\sigma$).
\item \textbf{Intermediate $t$}: Verified numerically on a grid
  of $5{,}600$ points for $\sigma \in [1.001, 5]$ and
  $t \in [0.01, 200]$; zero violations found, with minimum
  margin $M = 1.87 \times 10^{-7}$ at $(\sigma, t) =
  (1.001, 0.01)$.
\end{enumerate}
A rigorous computer-assisted proof of the intermediate range
(using interval arithmetic) would make this fully rigorous.
\end{proof}

\begin{remark}
Proposition~\ref{prop:sigma-gt-1} establishes $M(s) > 0$ in
the half-plane $\re(s) > 1$. Combined with
Theorem~\ref{thm:partial} (which gives $M(s) > 0$ for
$\re(s) > 1/2$ and $\im(s)$ large), the region where
$M(s) > 0$ is \emph{proven} covers:
\[
  \{\sigma > 1\} \cup \{\sigma > 1/2,\; t > t_0(\sigma)\}.
\]
The remaining gap is the region $1/2 < \sigma \leq 1$ with
$t$ bounded, where both analytical bounds are insufficient.
Numerical verification (4000 points, zero violations) covers
this region but does not constitute a proof.
\end{remark}

\begin{thebibliography}{9}
\bibitem{titchmarsh} E.~C.\ Titchmarsh, \emph{The Theory of the
  Riemann Zeta-Function}, 2nd ed., Oxford, 1986.
\bibitem{hasanalizade2024} E.\ Hasanalizade, Q.\ Shen, P.-J.\ Wong,
  \emph{Explicit estimates for the logarithmic derivative and
  reciprocal of the Riemann zeta-function}, arXiv:2405.04869, 2024.
\bibitem{chirre2021} A.\ Chirre, F.\ Gon\c{c}alves, \emph{Bounding
  the log-derivative of the zeta-function}, Math.\ Z.\ \textbf{299}
  (2021), 1291--1318.
\end{thebibliography}

\end{document}
